\subsection{可度量化拓扑向量空间中 0 的邻域}

我们称拓扑向量空间 E 是可度量化的，如果它的拓扑是可度量化的。相对于其加群的结构与其拓扑，E 因此是一个可度量化群（GT, IX, § 3.1）。

我们知道，为使一个拓扑群可度量化，必须且只需存在中性元 e 的一个可数的基本邻域系，它的交是单个元素 e（GT, IX, § 3.1, prop. 1）。

我们也知道，可度量化拓扑向量空间 E 的一致结构可以由一个不变距离 \(d(x, y) = |x - y|\) 定义，其中 \(x \mapsto |x|\) 是 E 到 \(R_{+}\) 中的一个连续映射，且满足条件：1) \(|-x| = |x|\)；2) \(|x + y| \leqslant |x| + |y|\)；3) 关系 \(|x| = 0\) 等价于 x = 0（GT, IX, § 3.1, prop. 3）。

我们已经看到（GT, IX, § 3.1, prop. 2），如何定义这样一个距离 \(d\)（利用 E 中 0 的邻域的一个递减序列 \((\mathbf{W}_n)\)），该序列构成一个基本邻域系，且满足 \(\mathbf{W}_{n+1} + \mathbf{W}_{n+1} + \mathbf{W}_{n+1} \subset \mathbf{W}_n\)。当 E 是非离散赋值除环 K 上的可度量化向量空间时，还可以假设诸 \(\mathbf{W}_n\) 是均衡的（I, 第 7 页, prop. 4）；回到 \(d\) 的定义过程（loc. cit.），我们可以看出关系 \(|\lambda| \leqslant 1\) 蕴含 \(|\lambda x| \leqslant |x|\)。此外，I, 第 2 页 的条件 (EVT \(_i\)) 与 (EVT \(_ii\)) 既蕴含 \(|\lambda x_0|\)（当 \(\lambda\) 在 K 中趋于 0 时，对每个 \(x_0 \in \mathrm{E}\)）趋于 0，又蕴含 \(|\lambda_0 x|\)（当 \(|x|\) 趋于 0 时，对每个 \(\lambda_0 \in \mathrm{K}\)）趋于 0。反之，若函数 \(|x|\) 具有前述所有性质，且 \(\mathbf{W}_n\) 是由 \(x \in \mathrm{E}\)（满足 \(|x| \leqslant 2^{-n}\)）所成的集合，则诸 \(\mathbf{W}_n\) 构成 E 上与 E 的向量空间结构相容的某个可度量化拓扑的 0 的均衡邻域基本系。

注记. --- 最重要的可度量化向量空间类之一是赋范空间（I, 第 3 页）。但必须注意，存在其拓扑不能用范数定义的可度量化向量空间（I, § 3, 习题 1）；我们稍后将研究一些重要的例子。

